El proyecto surgió como respuesta a las limitaciones operativas detectadas en un brazo robótico cartesiano desarrollado en un trabajo de graduación previo. El sistema existente presentaba deficiencias críticas en el desplazamiento de sus ejes, las cuales no pudieron ser solventadas mediante métodos de mantenimiento convencional, como la lubricación. Ante este escenario, el trabajo planteó una reingeniería de los componentes de movimiento mediante la sustitución de los ejes originales por ejes de acero y la integración de rodamientos lineales de bolas. Esta modificación tuvo como propósito reducir la fricción y los movimientos irregulares, de manera que se mejorara la fluidez cinemática para obtener un funcionamiento más continuo y preciso.

La relevancia del rediseño mecánico radicó en la necesidad de automatizar el brazo para su integración en una línea de empaque simulada dentro del Laboratorio de Diseño de Procesos de la Universidad del Valle de Guatemala. En este contexto, se consideró necesario que el sistema de transferencia mantuviera tiempos de ciclo reducidos y una repetibilidad constante para evitar interrupciones en el flujo del proceso. Mediante la automatización del brazo, se buscó reducir la dependencia del tiempo de reacción de un operador humano y disminuir la variabilidad asociada a la ejecución manual de tareas repetitivas. El proyecto se orientó al desarrollo de un sistema capaz de reaccionar ante la presencia de productos y ejecutar ciclos de manipulación de manera consistente, manteniendo condiciones similares de precisión y desempeño mecánico.

Finalmente, la implementación del brazo automatizado permitió disponer de una plataforma orientada a la capacitación y experimentación de estudiantes de ingeniería dentro de un entorno controlado. La integración de un sistema de control con sensores de posición y visión computacional permitió coordinar los movimientos del brazo con la información obtenida del proceso. De esta manera, el proyecto no solamente abordó las limitaciones mecánicas identificadas en el sistema original, sino que también contribuyó al desarrollo de una plataforma académica para el estudio de automatización, robótica, sistemas embebidos, sensores y visión computacional.